\documentclass[11pt]{article}
\usepackage[a4paper,margin=1in]{geometry}
\usepackage{booktabs}
\usepackage{graphicx}
\usepackage{listings}
\usepackage{amsmath}
\usepackage{hyperref}

\title{Problem 6: Parallel Optimization of the NEP Potential}
\author{}
\date{June 20, 2026}

\begin{document}
\maketitle

\section{Objective}
Problem 6 asks for NEP-specific parallel optimization, including analysis of
the calculation flow, optimization of atomic-density calculations,
multi-element processing, memory access, correctness testing, and performance
measurement over system sizes and thread counts.

This work extends the OpenMP force optimization from Problem 2 to the full NEP
CPU pipeline. It targets neighbor construction, descriptor/density evaluation,
mixed-element execution, and force calculation.

\section{NEP Calculation Flow and Hotspots}
For each molecular-dynamics step, \texttt{NEP::compute} performs:
\begin{enumerate}
  \item periodic-boundary processing and radial/angular neighbor-list construction;
  \item radial and angular atomic descriptor calculation;
  \item element-dependent neural-network evaluation for atomic energy and derivatives;
  \item radial, angular, and optional ZBL force/virial accumulation.
\end{enumerate}

The descriptor stage is an atomic-density calculation: every center atom
accumulates radial and angular contributions from its neighbors. Center atoms
are independent during descriptor construction, making atom-level OpenMP
parallelism natural. Force calculation is more difficult because pair
contributions update two atoms and require synchronization or reduction.

The remaining NEP-specific bottleneck found in this work was the large-box
cell-list neighbor search. Its expensive outer loop over center atoms was
serial even though each center atom writes to a separate neighbor-list column.

\section{Implementation}

\subsection{Large-Box Neighbor Parallelism}
The outer center-atom loop in \texttt{find\_neighbor\_list\_large\_box} now uses:
\begin{lstlisting}[language=C++]
#pragma omp parallel for schedule(static)
\end{lstlisting}
The temporary cell coordinate was moved inside the loop so every thread has a
private copy. Shared cell counts and cell contents are read-only in this phase,
while each thread writes only entries associated with its own center atom.
Periodic-boundary coordinate conversion is also parallelized over atoms.

This optimization is especially relevant for the 512-atom benchmark, whose
box is large enough to use the cell-list path.

\subsection{Descriptor and Atomic-Density Parallelism}
The standard and charge-aware NEP descriptor loops use explicit static
scheduling. Each thread owns one center atom at a time and stores descriptor,
network derivative, and angular density data at indices belonging to that atom.
Model parameters and element-dependent coefficient arrays are read-only.

Static scheduling was selected because the benchmark systems have comparable
local density per atom. It avoids the queue-management overhead of dynamic
scheduling and preserves predictable memory access.

\subsection{Multi-Element Processing}
Element type is read independently for every center-neighbor pair. The OpenMP
loops do not group or serialize element types, so O, Hf, and cross-element
interactions execute concurrently. Correctness and performance are tested for:
\begin{itemize}
  \item \texttt{single}: all atoms use element type 0;
  \item \texttt{mixed}: alternating type 0/type 1 atoms;
  \item \texttt{skewed}: a 3:1 ratio of type 0 to type 1.
\end{itemize}

\subsection{Force and Memory Strategy}
The optimized default force path uses thread-local force and virial buffers,
followed by one reduction per thread. This removes atomics from the inner
neighbor loop. NEP data remains in the existing structure-of-arrays layout:
separate $x$, $y$, and $z$ coordinate/force regions and neighbor-major buffers.
The implementation therefore avoids interface-changing data conversions.

Thread-local reduction costs $O(TN)$ extra memory, where $T$ is the thread
count and $N$ is the atom count. The Problem 2 serial and atomic force modes
remain available for validation and memory/performance comparison.

\section{Correctness Tests}
The NEP OpenMP unit test now includes a mixed-element descriptor test comparing
one and four threads. It also compares serial, atomic, and thread-local force
paths. The dedicated Problem 6 benchmark compares every descriptor, potential,
force, and virial array with a one-thread reference.

Across 27 Problem 6 benchmark cases, the maximum observed differences were:
\[
\max|\Delta D|=0,\qquad \max|\Delta E|=0,
\]
\[
\max|\Delta F|=7.10543\times10^{-15},\qquad
\max|\Delta V|=4.26326\times10^{-14}.
\]
Descriptor and energy results are identical. Force and virial differences are
normal floating-point reduction-order effects and are much smaller than the
unit-test tolerance of $10^{-10}$.

\section{Benchmark Script}
The dedicated files are:
\begin{lstlisting}
source/source_md/tools/nep_problem6_benchmark.cpp
source/source_md/tools/run_problem6_nep_benchmark.sh
\end{lstlisting}

The command used to generate the retained CSV was:
\begin{lstlisting}[language=bash]
cd source/source_md
tools/run_problem6_nep_benchmark.sh \
  ../../tests/PP_ORB/nep_hfo2.txt \
  64,216,512 1,2,4 single,mixed,skewed 10 \
  reports/problem6_nep_results.csv
\end{lstlisting}

The script compiles with \texttt{g++-15 -O2 -fopenmp}. Model loading is
excluded from timing. Descriptor timing includes neighbor construction and
atomic-density evaluation; compute timing covers the complete NEP step.

\section{Performance Results}
Table~\ref{tab:p6} reports four-thread speedup relative to one thread.

\begin{table}[ht]
\centering
\small
\resizebox{\textwidth}{!}{%
\begin{tabular}{@{}llrrrr@{}}
\toprule
Atoms & Composition & Descriptor 1T (ms) & Descriptor 4T speedup & Compute 1T (ms) & Compute 4T speedup \\
\midrule
64  & Single & 1.4195  & 3.38 & 3.2182  & 2.86 \\
64  & Mixed  & 1.2691  & 2.97 & 3.2470  & 2.94 \\
64  & Skewed & 1.2723  & 3.06 & 3.2246  & 2.96 \\
\midrule
216 & Single & 3.5527  & 3.42 & 10.0926 & 3.43 \\
216 & Mixed  & 3.4916  & 3.31 & 10.2138 & 3.52 \\
216 & Skewed & 3.5803  & 3.47 & 10.2689 & 3.53 \\
\midrule
512 & Single & 12.5353 & 3.68 & 28.4781 & 3.68 \\
512 & Mixed  & 12.7083 & 3.86 & 30.3105 & 4.01 \\
512 & Skewed & 12.4331 & 3.56 & 29.3249 & 3.84 \\
\bottomrule
\end{tabular}
}
\caption{Problem 6 NEP descriptor and full-compute scaling.}
\label{tab:p6}
\end{table}

The full raw table, including 1/2/4-thread times and all numerical differences,
is stored in \texttt{reports/problem6\_nep\_results.csv}.

\section{Analysis}
Descriptor/density calculation scales well for all compositions. Four-thread
speedup ranges from $2.97\times$ to $3.86\times$. The nearly identical scaling
for single, balanced, and skewed compositions demonstrates that element lookup
and cross-element coefficient selection are thread-safe and do not introduce a
serial multi-element bottleneck.

Full NEP compute reaches $2.86\times$--$4.01\times$ speedup. The 512-atom
systems exercise the newly parallelized large-box cell-list path and achieve
$3.68\times$--$4.01\times$ scaling. The mixed 512-atom case decreases from
30.3105 ms to 7.5596 ms.

The 512-atom single-element two-thread descriptor result is less efficient than
the surrounding cases, likely due to runtime scheduling and measurement noise;
its four-thread result returns to $3.68\times$ speedup. Repeated measurements
or CPU affinity control would improve benchmark stability on a shared desktop.

\section{Requirement Coverage and Conclusion}
The implementation addresses the core Problem 6 requirements:
\begin{itemize}
  \item the complete NEP flow and hotspots were analyzed;
  \item large-box neighbor and density preparation were parallelized;
  \item descriptor loops use explicit OpenMP scheduling;
  \item single- and multi-element systems were verified;
  \item memory behavior and force reductions were considered;
  \item correctness was checked across sizes, compositions, and threads;
  \item performance was measured for 27 configurations;
  \item a unit test, reproducible script, CSV, and report were produced.
\end{itemize}

The optimized thread-local NEP path is numerically consistent and demonstrates
strong four-thread scaling. Future work could add atom blocking to reduce
thread-local memory, adaptive cutoff selection, GPU density kernels, and
many-core server benchmarks.

\end{document}
